% block: shifted-systems-heading pages:175,176 type:section
\subsection{\BT{Shifted Linear Systems}{移位线性方程组}}\label{sec:shifted-linear-systems}
\BT{The Krylov subspace methods in Chap.\ \ref{chap:krylov-classification} are attractive for solving the following linear systems:}{Chapter\ \ref{chap:krylov-classification} 中的 Krylov 子空间方法对于求解下列线性方程组很有吸引力:}
\begin{equation}\label{eq:shifted-linear-systems}A^{(k)}\bm x^{(k)}=\bm b,\quad k=1,2,\ldots,m,\end{equation}
\BT{where $A^{(k)}$ is a nonsingular matrix of the form}{其中 $A^{(k)}$ 为下列形式的非奇异矩阵}
\begin{equation}\label{eq:shifted-coefficient-matrix}A^{(k)}:=A+\sigma_k I,\end{equation}
\BT{with $I$ being the identity matrix and $\sigma_k\in\mathbb C$. The linear systems \eqref{eq:shifted-linear-systems} are referred to as shifted linear systems. For simplicity, we also use the following notation:}{这里 $I$ 为单位矩阵, $\sigma_k\in\mathbb C$. 线性方程组 \eqref{eq:shifted-linear-systems} 称为移位线性方程组. 为简便起见, 也使用下列记号:}
\[(A+\sigma I)\bm x^{\sigma}=\bm b.\]
\BT{Here $\alpha^{\sigma},\beta^{\sigma},\ldots$ and $\bm a^{\sigma},\bm b^{\sigma},\ldots$ denote scalars and vectors related to the shifted linear system $(A+\sigma I)\bm x^{\sigma}=\bm b$. Note that $\alpha^{\sigma}$ is not the $\alpha$ to the power of $\sigma$.}{这里 $\alpha^{\sigma},\beta^{\sigma},\ldots$ 和 $\bm a^{\sigma},\bm b^{\sigma},\ldots$ 表示与移位线性方程组 $(A+\sigma I)\bm x^{\sigma}=\bm b$ 有关的标量和向量. 注意, $\alpha^{\sigma}$ 不是 $\alpha$ 的 $\sigma$ 次幂.}
% block: shifted-krylov-invariance page:176 type:ordinary-paragraph
\BT{If a Krylov subspace method with the initial guess $\bm x_0=\bm0$ is applied to \eqref{eq:shifted-linear-systems}, then from \eqref{eq:krylov-iterate-condition} the Krylov subspace method finds the approximate solutions of \eqref{eq:shifted-linear-systems} over the following subspaces:}{若以初始猜测 $\bm x_0=\bm0$ 将 Krylov 子空间方法应用于 \eqref{eq:shifted-linear-systems}, 则由 \eqref{eq:krylov-iterate-condition}, 该 Krylov 子空间方法在下列子空间上寻找 \eqref{eq:shifted-linear-systems} 的近似解:}
\begin{equation}\label{eq:shifted-krylov-membership}\bm x_n^{(k)}\in\mathcal K_n(A^{(k)},\bm b),\quad k=1,2,\ldots,m.\end{equation}
\BT{From the definition of Krylov subspace \eqref{eq:krylov-subspace-span}, the following shift invariance property holds true:}{由 Krylov 子空间的定义 \eqref{eq:krylov-subspace-span}, 下列移位不变性成立:}
\begin{equation}\label{eq:krylov-shift-invariance}\mathcal K_n(A,\bm b)=\mathcal K_n(A^{(k)},\bm b).\end{equation}
% block: shift-invariance-example page:176 type:ordinary-paragraph
\BT{This can be easily shown by induction and the definition of a Krylov subspace. As an example, for the case $n=3$ it follows that}{利用归纳法和 Krylov 子空间的定义很容易证明这一点. 例如, 当 $n=3$ 时, 有}
\[\begin{aligned}\mathcal K_3(A+\sigma I,\bm b)&=\{c_1\bm b+c_2(A+\sigma I)\bm b+c_3(A+\sigma I)^2\bm b:c_1,c_2,c_3\in\mathbb C\}\\
&=\{(c_1+c_2\sigma+c_3\sigma^2)\bm b+(c_2+2c_3\sigma)A\bm b+c_3A^2\bm b:c_1,c_2,c_3\in\mathbb C\}\\
&=\{d_1\bm b+d_2A\bm b+d_3A^2\bm b:d_1,d_2,d_3\in\mathbb C\}\\
&=\mathcal K_3(A,\bm b).\end{aligned}\]
\BT{In the third equation, the relation between $c_i$ and $d_i$ is given by}{在第三个等式中, $c_i$ 与 $d_i$ 的关系为}
\[\begin{bmatrix}d_1\\d_2\\d_3\end{bmatrix}=\begin{bmatrix}1&\sigma&\sigma^2\\0&1&2\sigma\\0&0&1\end{bmatrix}\begin{bmatrix}c_1\\c_2\\c_3\end{bmatrix}.\]
% block: collinear-shift-invariance page:176 type:ordinary-paragraph
\BT{Since the $3\times3$ matrix is nonsingular, for any $d_1,d_2,d_3$ there exist $c_1,c_2,c_3$. Thus the third equation holds true. From this example, we see that the shift invariance property can be slightly extended as follows:}{由于该 $3\times3$ 矩阵非奇异, 对任意 $d_1,d_2,d_3$ 都存在 $c_1,c_2,c_3$. 因而第三个等式成立. 由此例可见, 移位不变性可以略作推广如下:}
\[\mathcal K_n(A,\bm v)=\mathcal K_n(A^{(k)},\bm w),\quad\bm v=c\bm w.\]
\BT{If $\bm v=c\bm w$ holds for some $c$, the two vectors are referred to as collinear.}{若对某个 $c$ 有 $\bm v=c\bm w$, 则称这两个向量共线.}
% block: shared-shifted-subspace pages:176,177 type:ordinary-paragraph
\BT{Now from the shift invariance property \eqref{eq:krylov-shift-invariance}, it follows that \eqref{eq:shifted-krylov-membership} can be rewritten as}{现在由移位不变性 \eqref{eq:krylov-shift-invariance} 可知, \eqref{eq:shifted-krylov-membership} 可改写为}
\[\bm x_n^{(k)}\in\mathcal K_n(A,\bm b),\quad k=1,2,\ldots,m.\]
\BT{This means that only one Krylov subspace, $\mathcal K_n(A,\bm b)$, is required for obtaining the approximate solutions $\bm x_n^{(1)},\bm x_n^{(2)},\ldots,\bm x_n^{(m)}$. Note that if the shift-invariance property does not hold, $m$ different Krylov subspaces are required. In the following sections, we will see how to utilize the shift-invariance property and how this property is incorporated into Krylov subspace methods.}{这意味着只需一个 Krylov 子空间 $\mathcal K_n(A,\bm b)$, 就可得到近似解 $\bm x_n^{(1)},\bm x_n^{(2)},\ldots,\bm x_n^{(m)}$. 注意, 若移位不变性不成立, 则需要 $m$ 个不同的 Krylov 子空间. 下面各节将介绍如何利用移位不变性, 以及如何将这一性质融入 Krylov 子空间方法.}
